\par\Needspace{12\baselineskip}\pdfbookmark[1]{Vega. Thesaurus Logarithmorum. 1851.}{vega1851}
\begin{center}
\rule{0.7\linewidth}{0.4pt}\par
Astronomische Nachrichten. No.\ 756. 1851. May 2.\par
\rule{0.7\linewidth}{0.4pt}
\end{center}
\begin{center}\emph{Some Remarks on} \textsc{Vega}'s \emph{Thesaurus Logarithmorum, by}\par
\emph{Privy Councillor} \textsc{Gauss}.\end{center}

As is well known, \textsc{Vega}'s \emph{Thesaurus Logarithmorum} is for the most part a new printing of \textsc{Vlacq}'s larger logarithm tables. In the preface, \textsc{Vega} lists a considerable number of errors in the original that he has corrected, adding that he has also put right a very large number of inaccuracies in the last digit of the logarithms, amounting to one, two, three or four units. The newly added tables too, he says, were calculated, checked and corrected with equal care (in particular, therefore, the logarithms of the trigonometric quantities given for every individual second in \SourceLeaf{267} the first two degrees). \textsc{Vega} now seems to have flattered himself with the hope that his tables had in this way become almost free of errors. To arrive at \emph{completely} error-free tables, he promises a reward of one ducat for the first report of each error that might still remain and could lead to an incorrect calculation (\emph{pro sphalmatibus calculum turbantibus}). Whether this offer, dated $1^{\mathrm{st}}$ October $1794$, ever had any result is unknown to me.

I am unsure whether \textsc{Vega} had made entirely clear to himself which errors should be regarded as possibly leading to incorrect calculations. For all tables intended to represent theoretically determined irrational quantities, the familiar principle is that the tabulated quantity should always approach the true value as closely as the chosen number of decimal places permits; consequently, the deviation must never exceed half a unit in the last decimal place. Every breach of this strict rule is an error that could lead a calculator who relies on strict agreement to an incorrect result. If this strict interpretation is abandoned and the \emph{significance} of the inaccuracy is allowed to enter the judgement, the decision strays into the realm of arbitrariness. The circumstance already mentioned, that \textsc{Vega} himself speaks of corrections to \textsc{Vlacq}'s tables amounting to only one unit in the last place, and describes complete freedom from error as his goal, does indeed seem to indicate that he had the strict interpretation in mind. Moreover, in very many comparisons made as occasion arose with determinations to more places, I have always found the first part of the Thesaurus, which contains the logarithms of consecutive integers, very accurate.

Since then, similar prizes for the first report of errors in the numbers have been offered for several other logarithm tables, with the aim of perfecting their accuracy. I do not know, however, whether these offers had any success, apart from that concerning the logarithmic and trigonometric handbook edited by \textsc{Köhler} and published by Tauchnitz in Leipzig in $1847$. When these tables first appeared, the publisher promised one louis d'or for the first \SourceLeaf{268} report, submitted within a specified period, of each error capable of causing incorrect results. According to a printed report of 1 July 1848, this reward was actually paid for four errors reported. What deserves honourable mention here is that one of these errors could qualify only if the strict interpretation above was accepted. The logarithm of \(103000\), which to \(12\) places is
\[
=5{,}012837224705
\]
was given in the first printings with eight digits as
\[
=5{,}01283723
\]
whereas the number rounded according to the principle is \(5{,}01283722\). It is to be hoped that this decision will be respected as a precedent in similar cases in future.

Had \textsc{Vega} accepted the same strict interpretation, he might easily have fared like King Shiram, whose granaries did not suffice to pay the reward he had promised the inventor of chess.

\Needspace{42\baselineskip}
That the reward offered by \textsc{Vega} could, on that assumption, have cost him dearly can already be seen, without any recalculation, from a circumstance that is easily confirmed but, to my knowledge, has not been discussed anywhere else. In the table of logarithms of trigonometric quantities, the numbers in the sine column are almost without exception\HistoricalNote{21} exactly equal to the sum of the numbers in the cosine and tangent columns. Since all these numbers are only rounded values of the exact irrational quantities, it is clear that with strictly correct rounding this equality will fail in every case where the deviations from the exact values in the second and third columns have the same sign and their sum exceeds half a unit in the last decimal place. It is easy to see that among a large \SourceLeaf{269} number of cases this exception will occur, on average, once in four. Actual counts in seven-digit tables in which care has been taken over the correctness of the last digit confirm this. (For example, such a count in the \textsc{Köhler} tables mentioned, over the \(900\) individual minutes from \(30^{0}0'\) to \(44^{0}59'\), gave exactly \(225\) such cases; this perfect agreement must, of course, be regarded as accidental.) \textsc{Vega}'s Thesaurus contains the logarithms of the sines, cosines and tangents of \(22680\) arcs. On the assumption, therefore, that all the numbers in two of these columns have been rounded strictly according to the principle, one may expect, with a small uncertainty either way, \(5670\) logarithms in the third column to be incorrect by one unit. For the price of that many ducats, capable people would probably have been willing to recalculate the whole.

Which two columns are the original ones can be determined by comparison with other determinations reliable to more than ten places: from only a few arcs, and with certainty, if the assumption of strict correctness in the original columns holds; otherwise, from a somewhat larger number, at least with a preponderance of probability.

\Needspace{42\baselineskip}
Some material useful for this purpose is already available at the end of the decimal tables published by \textsc{Hobert} and \textsc{Ideler} in 1799. Although printed with only seven digits, they had been completed in manuscript to twice as many. There, \(138\) erroneous logarithms in \textsc{Vega}'s Thesaurus are listed. To conform exactly to the rule, \(127\) of them require a correction of one unit in the tenth place, \(10\) a correction of two units, and one a correction of three units. Of the \(127\), \(49\) concern sines, \(35\) cosines and \(43\) tangents; the remaining \(11\) corrections, of \(2\) or \(3\) units, concern tangents only. The arcs in question are those measured by multiples of \(27\) minutes; those not listed are the ones whose logarithms needed no correction. These comparisons could, moreover, have been made without \textsc{Hobert-Ideler}'s manuscript tables, since the widely available \textsc{Callet} tables contain the logarithms of sines and cosines to \(14\) digits for every thousandth of the quadrant. If one also takes from them the material needed for all arcs below \SourceLeaf{270} two degrees that are measured by multiples of \(5'24''\) but not of \(27'\), one obtains a comparison of \textsc{Vega}'s logarithms for \(118\) arcs with more accurate determinations. The results are assembled in the following summary (I):
\[
\begin{array}{rccc}
 & \mathrm{Sin.} & \mathrm{Cos.} & \mathrm{Tang.} \\
0 & 65 & 75 & 54 \\
1 & 53 & 43 & 53 \\
2 &    &    & 10 \\
3 &    &    & 1
\end{array}
\]
The meaning of these numbers, for example those in the last column, is that among \(118\) tangent logarithms, \(54\) are correctly given, \(53\) need a correction of one unit in the tenth place, \(10\) a correction of two units, and one a correction of three units. This already shows decisively that \emph{none} of the columns in \textsc{Vega}'s table is correct throughout, and with a preponderance of probability that the tangent logarithms are the derived ones, obtained by simply subtracting the numbers in the cosine column from those in the sine column.

\Needspace{22\baselineskip}
The authors of the decimal tables mentioned could, moreover, have supplied twice as many comparisons from their manuscript. A much larger yield can be obtained, however, by using the very valuable table of \textsc{Brigg}'s in the Trigonometria Britannica, published by \textsc{Gellibrand} in 1633 after its author's death. It gives the logarithms of sines and cosines to \(14\) digits for every hundredth of an ordinary degree, and \textsc{Callet} took the numbers mentioned above from it. It supplies the material for checking \textsc{Vega}'s tables at \(1060\) arcs, or \(3180\) logarithms altogether. I myself have confined this check to \(81\) arcs, or \(243\) logarithms, from \(14^{0}0'\) to \(18^{0}0'\). The result follows here (II):
\[
\begin{array}{rccc}
 & \mathrm{Sin.} & \mathrm{Cos.} & \mathrm{Tang.} \\
0 & 29 & 56 & 36 \\
1 & 52 & 25 & 42 \\
2 &    &    & 3
\end{array}
\]
\par
\SourceLeaf{271}
Combining groups (I) and (II), but counting the two arcs common to both only once, gives the following yield for 190 arcs (III):
\begin{center}
\begin{tabular}{rrrr}
 & Sin. & Cos. & Tang. \\
0 & 89 & 126 & 87 \\
1 & 101 & 64 & 90 \\
2 & & & 12 \\
3 & & & 1
\end{tabular}
\end{center}
Before these numbers can be used to estimate the ratio of correct to incorrect logarithms, a small deduction must be made. In (I), the logarithm of the sine of $45^{0}$ has been counted twice, namely also as $\log\cos45^{0}$; the tangent column also includes the logarithm of the tangent of $45^{0}$, which is rational. The same applies to (III). For the purpose above, its result must therefore be stated as follows: of 568 irrational logarithms checked, 301 proved correct and 267 incorrect. If this ratio could be regarded as holding on average, among the 68038 irrational logarithms in \textsc{Vega}'s Thesaurus, excluding, as is reasonable, the cotangents, probability would lead one to expect about 31983 errors.

\Needspace{36\baselineskip}
This number is probably still much too small. Among my papers I find calculations of the trigonometric logarithms to 14 digits, made some time ago for other purposes, for a couple of groups of arcs that advance, not in jumps, but at the same intervals as \textsc{Vega}'s tables. These show at least that the result (III) above does not yet give a true measure of the inaccuracy of these tables. Among those 190 arcs, there is not a single case in which the logarithm of a sine or cosine would need correction by more than \textit{one} unit; such cases are by no means rare among arcs that do \textit{not} occur in the Trigonometria Britannica, and for which the material needed for comparison must therefore first be obtained by a separate calculation. I therefore append the results of those calculations, since they can give our conception of the degree of inaccuracy in the numbers of \textsc{Vega}'s Thesaurus a firmer basis.
\par
\SourceLeaf{272}
Comparison of the numbers in the Thesaurus for the 21 arcs from $15^{0}38'20''$ to $15^{0}41'40''$ with more accurately calculated values gave (IV):
\begin{center}\begin{tabular}{r r r r}
& Sin. & Cos. & Tang.\\
0&4&12&1\\
1&9&8&8\\
2&6&1&6\\
3&2&&4\\
4&&&2
\end{tabular}\end{center}
The greatest deviations occur at the following arcs:
\[
\begin{array}{l l r r r}
15^{0}40'20'',&\text{where the corrections}&+3,&-1,&+4\\
15\quad41\quad30,&\text{where the following}&+3,&0,&+4
\end{array}
\]
are to be applied to the logarithms of the sine, cosine and tangent.

\Needspace{11\baselineskip}Similarly, comparison of \textsc{Vega}'s numbers for the 93 arcs from $1^{0}19'52''$ to $1^{0}21'24''$ with a more accurate calculation gave the following (V):
\begin{center}\begin{tabular}{r r r r}
& Sin. & Cos. & Tang.\\
0&38&30&17\\
1&39&56&41\\
2&15&7&22\\
3&1&&11\\
4&&&2
\end{tabular}\end{center}
The greatest deviations occur at the arcs $1^{0}20'10''$ and $1^{0}20'15''$; the necessary corrections are $-2,+1,-4$ for the former and $-2,+2,-4$ for the latter.

\Needspace{28\baselineskip}
Groups IV and V do not have the same origin. The numbers in the Thesaurus for group IV were taken from \textsc{Vlacq}'s Trigonometria Artificialis, whereas the others belong to the newly added material calculated under \textsc{Vega}'s direction and supervision by Lieutenant \textsc{Dorfmund}. In the summaries above, the latter appear somewhat less inaccurate than the former, although group IV contains too few cases to support a firm judgement. In any event, combining the two groups gives an estimate \SourceLeaf{273} of the total inaccuracy of the tables that is more likely to be somewhat too favourable than the reverse. The combination gives, for 114 arcs (VI):
\begin{center}\begin{tabular}{r r r r}
& Sin. & Cos. & Tang.\\
0&42&42&18\\
1&48&64&49\\
2&21&8&28\\
3&3&&15\\
4&&&4
\end{tabular}\end{center}
Thus, among 342 logarithms, there are only 102. that need no correction, as against 240 inaccurate ones. At this ratio, among the 68038 irrational logarithms one would expect 47746 inaccurate ones.

The sum of the squares of the deviations is found to be 159 for sines, 96 for cosines and 360 for tangents. The mean error may therefore be taken as $1{,}18$ for sines, $0{,}92$ for cosines and $1{,}78$ for tangents. After this, there can be no doubt that the tangent logarithms are the derived ones.

\begin{samepage}
The smaller inaccuracy of the numbers in the cosine column than in the sine column can surely be attributed, at least in part, to the simpler interpolation methods needed to fill the gaps in the former. Other causes may possibly have contributed, but only uncertain conjectures could be made about them.
\begin{center}\bigskip\rule{.3\linewidth}{.4pt}\end{center}
\end{samepage}
